% Capítulo 2: Levantamento Bibliográfico (Estado da Arte)
\chapter{Levantamento Bibliográfico}\label{chp:levantamento}

Este capítulo apresenta a revisão da literatura sobre imageamento
hiperespectral, abordando os fundamentos, tecnologias disponíveis e
aplicações em sistemas embarcados. A análise foi organizada em sete
seções principais: princípios físicos da espectroscopia de imageamento
e formação de imagens espectrais; câmeras hiperespectrais, incluindo
tipos, modelos e datasets públicos; técnicas de processamento de
imagens hiperespectrais; tipos de processadores e suas características
para processamento hiperespectral; arquiteturas heterogêneas de
processamento; plataformas específicas para sensoriamento remoto; e
benchmarking de plataformas comerciais. A revisão abrange trabalhos
fundamentais da área e publicações recentes sobre aplicações em
agricultura de precisão, monitoramento ambiental e sistemas embarcados,
com ênfase nos desafios computacionais do processamento de grandes
volumes de dados e soluções para operação em tempo real com
restrições energéticas.

\section{Princípios Físicos da Espectroscopia de Imageamento e
    Formação de Imagens Espectrais}
\label{sec:fundamentos_espectro}

Os conceitos fundamentais do imageamento hiperespectral baseiam-se nos
princípios da física óptica que descrevem a interação entre a radiação
eletromagnética e a matéria \cite{Lorenzzetti2015}. Esta seção examina
as características das ondas eletromagnéticas e sua organização
espectral, com foco particular na faixa de 400 a 2500 nm, que abrange
as regiões do visível, infravermelho próximo e infravermelho de ondas
curtas, essenciais para o sensoriamento remoto. As propriedades ópticas
dos materiais são analisadas através dos fenômenos de reflectância,
transmitância e absorção, processos físicos que originam as assinaturas
espectrais únicas de cada substância e possibilitam sua identificação e
caracterização \cite{Baptista2019}.

A compreensão da distinção entre sistemas multiespectrais e
hiperespectrais, quanto à resolução espectral e capacidade de
discriminação, fundamenta a escolha da tecnologia adequada para cada
aplicação. Esses princípios teóricos são demonstrados através de
aplicações práticas em agricultura de precisão e monitoramento
ambiental, evidenciando como o conhecimento dos fundamentos físicos é
indispensável para o desenvolvimento de metodologias robustas de
aquisição, processamento e análise de dados hiperespectrais.

\begin{figure}[htbp]
    \centering
    \includegraphics[width=\textwidth]{conceito_espectrocospia_imageamento_hiperspectral}
    \caption{Conceito da espectroscopia de imageamento desenvolvido pelo sensor 
             AVIRIS (JPL/NASA). A figura ilustra como 224 bandas espectrais são 
             obtidas simultaneamente, permitindo a análise de diferentes materiais 
             (atmosfera, solos, água, vegetação) através de suas assinaturas espectrais 
             características. Este princípio fundamenta o processamento hiperespectral 
             em sistemas embarcados. Fonte: \cite{Baptista2019}.}
    \label{fig:conceito-espectroscopia-imageamento}
\end{figure}

\subsection{O espectro eletromagnético}
\label{subsec:espectro_eletromagnetico}

A energia eletromagnética constitui a base física do sensoriamento remoto
e do imageamento hiperespectral, podendo ter origem em fontes naturais,
como o Sol, nos próprios alvos ou em sistemas ativos, como radares. Essa
energia propaga-se em forma de ondas, transportando informações sobre os
materiais com os quais interage. A \autoref{fig:conceito-espectroscopia-imageamento} 
ilustra como o sensor AVIRIS implementa estes conceitos na prática, 
obtendo simultaneamente centenas de bandas espectrais. Compreender as 
propriedades fundamentais da radiação eletromagnética é essencial para 
entender os processos de aquisição e análise de dados hiperespectrais.

\subsubsection{Definição e características das ondas eletromagnéticas}

Ondas eletromagnéticas são perturbações que se propagam no espaço à
velocidade da luz ($c =3 \times 10^8$ m/s no vácuo). São
caracterizadas principalmente pelo comprimento de onda ($\lambda$) e
pela frequência ($f$), relacionados por:
\begin{equation}
c = \lambda \times f
\end{equation}
A energia ($E$) transportada por uma onda eletromagnética é proporcional
à sua frequência, conforme a equação de Planck:
\begin{equation}
E = h \times f = \frac{hc}{\lambda}
\end{equation}
em que $h$ é a constante de Planck ($6{,}626 \times 10^{-34}$~J$\times$s).

No contexto do sensoriamento remoto, diferentes unidades são empregadas para
descrever as faixas espectrais:
\begin{itemize}
    \item 1~kHz = $1 \times 10^{3}$~Hz
    \item 1~MHz = $1 \times 10^{6}$~Hz
    \item 1~GHz = $1 \times 10^{9}$~Hz
    \item 1~THz = $1 \times 10^{12}$~Hz
    \item 1~$\mu$m = $1 \times 10^{-6}$~m
    \item 1~nm = $1 \times 10^{-9}$~m
\end{itemize}

Neste trabalho adota-se a unidade de micrometro ($\mu$m) para o
comprimento de onda.

\subsubsection{Regiões espectrais: do ultravioleta ao infravermelho}

O espectro eletromagnético é dividido em regiões conforme as
propriedades físicas da radiação e sua interação com a matéria
\cite{Lorenzzetti2015, Baptista2019}. Para sensoriamento remoto,
destacam-se as seguintes faixas:

\begin{description}
    \item[Ultravioleta (UV):] 0,03 a 0,4~$\mu$m. Pouco utilizada
    em sensoriamento remoto terrestre devido à absorção atmosférica,
    mas relevante em aplicações específicas.
    \item[Visível:] 0,4 a 0,7~$\mu$m, correspondente à radiação
    percebida pelo olho humano. Subdivisões usuais:
\begin{itemize}
        \item Azul: 0,4---0,5~$\mu$m
        \item Verde: 0,5---0,6~$\mu$m
        \item Vermelho: 0,6---0,7~$\mu$m
\end{itemize}
    \item[Infravermelho próximo (NIR):] 0,7 a 1,3~$\mu$m. Importante
    para estudos de vegetação devido à alta reflectância da clorofila.
    \item[Infravermelho de ondas curtas (SWIR):] 1,3 a 3,0~$\mu$m.
    Fundamental para identificação de minerais e análise de umidade.
\end{description}

Os limites exatos dessas faixas podem variar entre autores, mas há consenso
sobre as principais subdivisões. A Tabela~\ref{tab:faixas_espectrais_iv}
apresenta as subfaixas do infravermelho mais utilizadas em sensoriamento
remoto.

\begin{table}[htbp]
    \centering
    \caption{Faixas espectrais e subdivisões no infravermelho utilizadas em
    sensoriamento remoto (valores em nanômetros, nm, até 1~mm).}
    \label{tab:faixas_espectrais_iv}
    \begin{tabular}{lcc}
    \toprule
    \textbf{Faixa} & \textbf{Subfaixa} & \textbf{Comprimento de onda} \\
    \midrule
    Raios-X & --- & 0{,}03---30~nm \\
    Ultravioleta & --- & 30---400~nm \\
    Visível & --- & 400---700~nm \\
    Infravermelho & NIR/IVP & 740---1\,000~nm \\
    & SWIR/IVOC & 1\,000---3\,000~nm \\
    & LWIR/IVTM & 3\,000---14\,000~nm \\
    Micro-ondas & --- & 1---30~mm \\
    Ondas de rádio & --- & >30~mm \\
    \bottomrule
    \end{tabular}
\end{table}

Observa-se que, em algumas referências \cite{Lorenzzetti2015, Baptista2019}, 
a faixa de micro-ondas pode ser
definida como indo de 1~mm até 1~m. Essas variações não comprometem a
aplicação prática dos conceitos, mas devem ser consideradas ao comparar
diferentes trabalhos.

\subsubsection{Comprimentos de onda relevantes para sensoriamento remoto
(400--2500 nm)}

Para aplicações de imageamento hiperespectral e sensoriamento remoto, a faixa
espectral de 400 a 2500 nm (0,4 a 2,5 $\mu$m) representa a região de maior
interesse prático. Essa faixa abrange as regiões do visível, infravermelho
próximo (NIR) e infravermelho de ondas curtas (SWIR), onde ocorrem os
principais fenômenos de interação radiação-matéria relevantes para
identificação e caracterização de materiais terrestres.

\textbf{Região VNIR (400--1000 nm):} A faixa do visível e infravermelho
próximo é fundamental para estudos de vegetação e classificação de cobertura
terrestre. Nessa região, processos eletrônicos dominam a absorção da radiação,
resultando em transições entre níveis de energia eletrônica nos átomos e
moléculas. A clorofila, por exemplo, apresenta forte absorção nas regiões azul
(430--450 nm) e vermelha (640--660 nm), enquanto exibe alta reflectância no NIR
(700--1000 nm).

\textbf{Região SWIR (1000--2500 nm):} No infravermelho de ondas curtas, as
absorções são governadas por processos vibracionais moleculares. Muitas das
feições encontradas nos espectros obtidos nas faixas do visível e do
infravermelho próximo decorrem de efeitos do campo cristalino, enquanto outros
processos eletrônicos produzem feições espectrais significativas nessa faixa.
Essa região é particularmente rica em informações sobre minerais, conteúdo de
água e compostos orgânicos.

As transições entre níveis de energia são responsáveis pelas faixas de absorção
nos espectros. As absorções são governadas por processos vibracionais, sendo
que átomos isolados e íons têm estados de energia discretos, caracterizando
eventos quânticos. A absorção de fótons em um comprimento de onda específico
resulta da mudança de um estado de energia mais baixo para um mais alto,
enquanto a emissão resulta da mudança de um estado de energia mais alto para um
mais baixo, com liberação de energia no mesmo comprimento de onda.

A compreensão dos fundamentos físicos da interação radiação-matéria é essencial
para o desenvolvimento de algoritmos de processamento hiperespectral eficazes e
para a interpretação adequada dos dados espectrais em aplicações práticas de
sensoriamento remoto.


\subsection{Propriedades Ópticas dos Materiais}
\label{subsec:propriedades_opticas}

As propriedades ópticas dos materiais constituem a base física que permite 
a identificação e caracterização de diferentes substâncias através do 
sensoriamento remoto hiperespectral \cite{Baptista2019}. Quando a 
radiação eletromagnética interage com um material, três processos 
fundamentais podem ocorrer: reflexão, absorção e transmissão. A compreensão 
destes processos e dos fatores que os influenciam é essencial para o 
desenvolvimento de técnicas eficazes de análise espectral e para a 
interpretação adequada dos dados hiperespectrais em aplicações práticas.

\subsubsection{Reflectância Espectral: Definição e Fatores Influenciadores}

A espectroscopia de reflectância é o estudo da radiação eletromagnética 
como função do comprimento de onda em que ela está sendo refletida, emitida 
ou espalhada por um gás, um líquido ou um sólido \cite{Baptista2019}. 
A energia refletida pelas superfícies pode ser entendida como a porção da 
energia incidente, subtraindo-se as perdas que foram absorvidas ou 
transmitidas.

A reflectância espectral $\rho(\lambda)$, função do comprimento de onda 
($\lambda$), pode ser definida matematicamente como a razão entre a 
radiância da superfície [L ($\lambda$)] expressa em W m$^{-2}$ sr$^{-1}$ 
$\mu$m$^{-1}$ no comprimento de onda ($\lambda$) e a irradiância incidente 
[E ($\lambda$)] em W m$^{-2}$ $\mu$m$^{-1}$ no comprimento de onda 
($\lambda$), conforme a equação:

\begin{equation}
\rho(\lambda) = \pi \frac{L(\lambda)}{E(\lambda)}
\end{equation}

Esta definição estabelece a base quantitativa para a medição e comparação 
das propriedades espectrais de diferentes materiais. 
O Sol é a principal fonte natural de energia eletromagnética nas faixas do 
espectro do ultravioleta, do visível, do infravermelho próximo e do 
infravermelho de ondas curtas. Em comprimentos de onda maiores, na faixa do 
infravermelho termal, a eficiência do Sol como fonte de energia é bastante 
reduzida e a Terra passa a ter maior eficiência como fonte, porém nesta 
faixa não se analisa o espectro óptico de reflectância e sim o de
emitância \cite{Baptista2019}.

O processo de interação dos fótons com a matéria resulta em que parte deles 
é refletida pelos grãos da superfície do material, alguns passam pelos 
grãos e alguns são absorvidos. Os fótons são 
absorvidos pelos minerais por diversos processos que dependem de sua 
composição química. Essas absorções são governadas por dois processos 
gerais: processos eletrônicos e processos vibracionais \cite{Baptista2019}.

\subsection{Princípios de Sensoriamento Remoto Hiperespectral}
\label{subsec:principios_sensoriamento}

\subsubsection{Interação Radiação-Matéria no Contexto de Sensoriamento Remoto}

O sensoriamento remoto hiperespectral fundamenta-se nos princípios físicos 
da interação entre radiação eletromagnética e matéria, constituindo uma 
extensão natural dos conceitos fundamentais abordados nas seções anteriores. 
A definição de sensoriamento remoto (SR) está sujeita a diferentes 
interpretações na literatura científica. De acordo com \textit{Manual of 
Remote Sensing} (Ulaby, 1983), bem como em Elachi e Zyl (2006), SR é 
definido como a aquisição de informação sobre alguma propriedade de um 
objeto ou fenômeno sem contato físico com ele. A informação sobre um alvo 
seria obtida pela detecção e medida de mudanças que o objeto impõe sobre o 
meio circundante, seja ele eletromagnético, acústico ou potencial.

Schowengerdt (2007) considera SR como ``a medida das propriedades de um 
objeto ou da superfície da Terra usando dados adquiridos de sensores e 
sistemas'', sendo um pouco mais restritivo. Segundo Slater (1980), SR é o 
conjunto de atividades utilizadas para a aquisição de informações relativas 
aos recursos naturais da Terra, ou ao seu meio ambiente, obtidas pela 
análise da energia eletromagnética refletida, emitida ou retroespalhada 
pelos alvos, coletada por meio de sensores remotos \cite{Baptista2019}.

No contexto hiperespectral, a interação radiação-matéria adquire uma 
dimensão adicional de complexidade devido à alta resolução espectral dos 
sistemas de imageamento. Enquanto os processos eletrônicos e vibracionais 
descritos anteriormente permanecem fundamentalmente os mesmos, a capacidade 
de detectar e quantificar feições espectrais sutis é significativamente 
ampliada. (CLARK, 1999 apud LORENZZETTI, 2015) considera que quatro 
parâmetros gerais definem a capacidade de um espectrorradiômetro (sensor 
espectral): a faixa espectral, a largura das bandas espectrais, a 
amostragem espectral e a razão sinal/ruído.

A faixa espectral de resolução utilizada em sensoriamento remoto abrange 
tipicamente os seguintes intervalos espectrais: de 0,3 a 1,0 $\mu$m, 
chamado de visível e infravermelho próximo; de 1,0 a 2,5 $\mu$m, 
infravermelho de ondas curtas. A faixa do infravermelho termal situa-se a 
partir de 3,0 $\mu$m até aproximadamente o pico máximo de emitância da 
Terra em torno de 10,0 $\mu$m \cite{Lorenzzetti2015}.

A largura da banda espectral é definida por (CLARK, 1999 apud LORENZZETTI, 
2015) como a largura de um canal espectral individual no 
espectrorradiômetro. Quanto mais estreita for a largura da banda espectral, 
mais precisa será a determinação da feição de absorção do mineral. Alguns sistemas sensores possuem bandas largas e não espectralmente
contíguas, não permitindo análise contínua do espectro, enquanto outros
sistemas, como o AVIRIS, apresentam largura de bandas espectrais estreitas
com espaçamento contínuo e permitem a identificação de feições
características dos minerais.

\subsubsection{Sensores Multiespectrais vs. Hiperespectrais: Diferenças 
Fundamentais}

A distinção fundamental entre sensores multiespectrais e hiperespectrais 
reside na resolução espectral e na capacidade de discriminação de materiais. 
O perfil da largura também é importante. O ideal é que cada canal de um 
espectrorradiômetro rejeite toda radiação que não esteja dentro do alcance 
estreito de determinado comprimento de onda. Mas o que geralmente ocorre, 
em virtude de efeitos ópticos muito complexos, é que a radiação pode vazar, 
difundindo-se dentro do sistema óptico ou dentro do sistema de filtros 
bloqueadores.

Normalmente a largura do perfil é definida como a largura em comprimento 
de onda em 50\% do nível de resposta da função do detector, chamado de 
largura a meia altura (FWHM) (CLARK, 1999 apud LORENZZETTI, 2015). A razão 
sinal/ruído (SNR) visa a obtenção de espectros com precisão suficiente para 
registrar todas as feições espectrais importantes.

A SNR é dependente da sensibilidade dos detectores, da largura da banda 
espectral da luz refletida ou emitida pela superfície. Algumas características espectrais são bastante fortes e uma razão
SNR baixa é suficiente para identificá-las, enquanto outras
feições são muito fracas e é necessária uma razão SNR mais alta (SWAYZE
et al., 1997 \textit{apud} LORENZZETTI, 2015).

Sistemas multiespectrais convencionais, como o sensor Landsat TM5, operam
com bandas espectrais relativamente largas e espaçadas, tipicamente entre
6 a 12 bandas cobrindo o espectro eletromagnético de interesse. Sistemas
como o Landsat e o \textit{Moderate Resolution Imaging Spectroradiometer}
(MODIS) não permitem a percepção de feições de absorção estreitas. Como exemplo, GREEN (1997 \textit{apud} LORENZZETTI, 2015) demonstra
como o sensor Landsat TM5 captaria as feições de absorção de três minerais:
hematita, caulinita e dolomita.

A caulinita apresenta sua feição espectral mais característica em 2,2 $\mu$m. 
Já a dolomita, assim como a calcita, apresenta essa feição próxima para sua 
identificação, em 2,3 $\mu$m. As feições espectrais utilizadas para 
identificação da caulinita e dolomita (ou da calcita, pois essas são 
semelhantes) não podem ser discriminadas, pois aparecem integradas na banda 
7 (2,08 a 2,35 $\mu$m) do sistema sensor TM5 do Landsat.

\subsubsection{Aplicações Práticas: Agricultura, Monitoramento 
Ambiental e Defesa}

As aplicações práticas do sensoriamento remoto hiperespectral em agricultura 
de precisão e monitoramento ambiental demonstram o potencial da tecnologia 
para resolver problemas reais com alta precisão e especificidade, 
constituindo um dos principais motivos para o desenvolvimento de sistemas 
embarcados eficientes.

\textbf{Agricultura de Precisão:}

Na agricultura de precisão, os sistemas hiperespectrais permitem o 
monitoramento detalhado da saúde vegetal, detecção precoce de estresses 
bióticos e abióticos, e estimativa de parâmetros biofísicos das culturas. 
A região do infravermelho próximo (700--1000 nm) é particularmente importante 
devido à alta reflectância da clorofila, permitindo a avaliação do vigor 
vegetativo através de índices espectrais específicos.

As feições de absorção no SWIR (1400 nm e 1900 nm) são diagnósticas para o 
conteúdo de água nas plantas, facilitando o monitoramento de estresse hídrico 
em tempo real. Esta capacidade é fundamental para sistemas embarcados em 
UAVs que operam em agricultura de precisão, permitindo tomada de decisões 
imediatas sobre irrigação e manejo de culturas.

A capacidade de discriminar espécies vegetais e estágios fenológicos através 
de suas assinaturas espectrais únicas representa um avanço significativo 
sobre os métodos convencionais. Sistemas hiperespectrais embarcados podem 
detectar variações sutis na composição bioquímica foliar, incluindo 
concentrações de clorofila, carotenoides, antocianinas e compostos 
estruturais como celulose e lignina.

\textbf{Monitoramento Ambiental:}

No monitoramento ambiental, a tecnologia hiperespectral oferece capacidades 
sem precedentes para a caracterização de materiais superficiais, qualidade 
da água e detecção de poluentes. A identificação de minerais específicos 
através de suas feições espectrais características permite o mapeamento 
geológico detalhado e a detecção de alterações hidrotermais associadas a 
depósitos minerais.

Para qualidade da água, os sensores hiperespectrais podem detectar e 
quantificar constituintes opticamente ativos como clorofila-a, sedimentos 
em suspensão, matéria orgânica dissolvida e diversos poluentes. As absorções 
características desses materiais em comprimentos de onda específicos 
permitem sua identificação e quantificação com precisão superior aos métodos 
convencionais.

A detecção de mudanças ambientais, sejam elas naturais ou antropogênicas, 
beneficia-se da alta sensibilidade espectral dos sistemas hiperespectrais. 
Alterações sutis na composição química ou física dos materiais superficiais, 
que seriam imperceptíveis em imagens multiespectrais, tornam-se detectáveis 
através da análise das variações nas assinaturas espectrais.

A integração desses princípios fundamentais com as tecnologias de 
processamento embarcado, tema central desta dissertação, representa o 
próximo passo evolutivo na aplicação prática do sensoriamento remoto 
hiperespectral, permitindo análises em tempo real e tomada de decisões 
imediatas em campo.

\textbf{Aplicações Militares e de Defesa:}

Grande parte das aplicações de sensoriamento remoto para fins militares 
é focada na detecção e classificação de alvos e no mapeamento de 
terrenos e instalações. Em geral, os dados obtidos (coletados por 
satélites ou por dispositivos aerotransportados) por aviões usados para 
aplicações militares são considerados como \textit{classificados} (não 
disponíveis para terceiros) e de alta resolução espacial. Como definido por Simonett et al. (1983) e Jensen (2009), considera-se
como resolução espacial de um sistema sensor a menor separação espacial
angular ou linear entre dois objetos que podem ser detectados.

Um detalhamento mais aprofundado sobre o conceito de resolução pode ser 
encontrado em Forshaw et al. (1983) e Townshend (1980). Como uma imagem 
digital é constituída por \textit{pixels}, costuma-se associar o tamanho 
do \textit{pixel} para a resolução da imagem. Entretanto, essa definição 
dá a impressão de que somente podemos resolver objetos de tamanho maior, 
ou igual ao tamanho do \textit{pixel}, o que nem sempre é verdade. Exemplos de definição mais rigorosa, considerando alto contraste com
pixels vizinhos, são apresentados por Schowengerdt (2007).

\textbf{Desafios Técnicos e Limitações:}

Os sistemas hiperespectrais enfrentam desafios únicos relacionados ao 
processamento de grandes volumes de dados em tempo real. A razão 
sinal/ruído (SNR) visa à obtenção de espectros com precisão suficiente 
para registrar todas as feições espectrais importantes. A SNR é dependente 
da sensibilidade dos detectores, da largura da banda espectral da luz 
refletida ou emitida pela superfície.

O teorema de Nyquist determina que o máximo de informação é obtido
quando a amostragem espectral é igual à metade do FWHM; às vezes, porém,
uma razão de amostragem diferente pode ser necessária. Muitos dos espectrômetros modernos em 
uso (por exemplo, o AVIRIS) possuem um intervalo de amostragem de 
aproximadamente 10 nm e FWHM de 10 nm (CLARK, 1999 \textit{apud} 
LORENZZETTI, 2015).

Algumas características espectrais são bastante fortes e uma razão
sinal/ruído (SNR) baixa é suficiente para identificá-las, enquanto outras
feições são muito fracas e é necessária uma razão SNR mais alta (SWAYZE
et al., 1997 \textit{apud} LORENZZETTI, 2015). Esta limitação é particularmente relevante para sistemas
embarcados, onde restrições de peso, energia e processamento computacional
impõem desafios adicionais à manutenção de alta qualidade espectral.

\section{Técnicas de Processamento de Imagens Hiperespectrais}
\label{sec:tecnicas_processamento_hiperespectral}

O processamento de imagens hiperespectrais envolve uma cadeia complexa de 
algoritmos especializados, desde o pré-processamento básico até técnicas 
avançadas de análise espectral. Esta seção apresenta as principais abordagens 
utilizadas na literatura, com foco nas técnicas que demonstram potencial para 
implementação em sistemas embarcados, consolidando metodologias validadas em 
diferentes contextos de aplicação.

\subsection{Pré-processamento e Correção Radiométrica}
\label{subsec:pre_processamento}

O pré-processamento constitui uma etapa crucial para assegurar a
qualidade dos dados hiperespectrais, especialmente em sistemas embarcados
que operam sob condições de iluminação e movimento variáveis. Por meio do
pré-processamento, os dados são corrigidos, permitindo comparações
consistentes entre diferentes séries temporais e até mesmo dados
provenientes de distintos sensores hiperespectrais. Em sistemas embarcados
em UAVs para a detecção de plásticos em ambientes aquáticos, o
pré-processamento demonstra ser fundamental para manter uma precisão
superior a 90\% em condições operacionais desafiadoras \cite{Alboody2023}.

As técnicas de pré-processamento mais importantes incluem a calibração
radiométrica, que corrige as variações de resposta do sensor utilizando
ferramentas de software fornecidas pelos fabricantes. Essa etapa é
essencial para georreferenciar e converter os cubos de dados DN (Digital
Number) brutos em radiância, garantindo a reprodutibilidade e
comparabilidade dos dados. Métodos avançados como o Empirical Line Method
(ELM) podem melhorar a precisão radiométrica entre 5\% e 55\%
\cite{SilvaSantos2024}.

A correção atmosférica remove os efeitos de absorção e espalhamento
atmosférico, sendo particularmente crítica em aplicações aéreas, com
algoritmos como ATCOR e FLAASH amplamente validados para essa finalidade.
A correção geométrica compensa as distorções geométricas causadas pelo
movimento da plataforma e pela geometria de aquisição, frequentemente
combinando dados de orientação externa coletados durante o voo com um
modelo do sensor e um modelo de elevação. Técnicas robustas de correção
geométrica são essenciais para aplicações embarcadas em UAVs, onde
variações de altitude e orientação introduzem distorções significativas
\cite{SilvaSantos2024}.

A redução de ruído espectral, através da aplicação de filtros adaptativos
e técnicas de suavização espectral, como os filtros Savitzky-Golay e
transformadas wavelets, demonstra eficácia na preservação de
características espectrais importantes, enquanto reduz componentes ruidosas.
Por fim, a avaliação da qualidade, onde a relação sinal-ruído (SNR)
constitui um critério fundamental para caracterizar a qualidade dos dados
hiperespectrais, sendo essencial para análises quantitativas precisas.
Sistemas comerciais como Specim FX10 e Resonon Pika L oferecem alta
qualidade espectral, adequada para aplicações exigentes de classificação
\cite{Zhang2024}.

\subsection{Redução Dimensional e Seleção de Bandas}
\label{subsec:reducao_dimensional}

A alta dimensionalidade dos dados hiperespectrais, com suas centenas de
bandas espectrais, impõe desafios computacionais significativos, tornando
as técnicas de redução dimensional essenciais para sistemas embarcados. O
fenômeno de Hughes demonstra como o desempenho de classificadores pode
degragar com o aumento da dimensionalidade quando o número de amostras de
treinamento é limitado, tornando essas técnicas ainda mais críticas.
Implementações otimizadas em FPGA de algoritmos como EMCR (Extreme Learning
Machine with Correlation Removal) conseguem reduzir 80\% dos dados,
mantendo uma precisão superior a 99\% \cite{Martins2019}.

Dentro das técnicas lineares de redução dimensional, destaca-se a Principal
Component Analysis (PCA), uma técnica clássica que projeta os dados em um
subespaço de menor dimensionalidade, preservando a máxima variância. A PCA
concentra a variância espectral e pode ser combinada com outros métodos, como
compressão wavelet, conseguindo retenção de 99\% da variância original com
redução de 80\% a 90\% das bandas. Implementações GPU de PCA para compressão
hiperespectral conseguem processamento em tempo real em plataformas embarcadas
como Jetson TX2 \cite{Diaz2019}.

O Maximum Noise Fraction (MNF) é outro método baseado em PCA que maximiza a
relação sinal-ruído, sendo capaz de alcançar boa performance na redução
dimensional de dados hiperespectrais, especialmente eficaz para dados com alta
correlação espectral. Aplicações em sistemas PRISMA demonstram eficácia superior
a 87\% na redução de ruído espectral \cite{Corbari2023}. A Independent Component
Analysis (ICA) decompõe os dados em componentes estatisticamente independentes,
sendo adequada para a separação de assinaturas espectrais sobrepostas e
identificação de endmembers puros.

Em relação aos métodos de seleção de bandas, a seleção baseada em entropia, com
algoritmos como EMCR (Entropy Multiple Correlation Ratio), consegue reduzir 70\%
a 80\% do processamento, mantendo alta precisão de classificação, sendo
computacionalmente eficientes para sistemas embarcados. Algoritmos genéticos
oferecem otimização global para seleção de subconjuntos ótimos de bandas, com
validação cruzada para evitar overfitting. Métodos híbridos combinam técnicas de
projeção com seleção de bandas, oferecendo um equilíbrio entre redução
dimensional e preservação de informações discriminantes.

\subsection{Algoritmos de Classificação Espectral}
\label{subsec:classificacao_espectral}

A classificação espectral constitui uma das aplicações mais importantes do
processamento hiperespectral, sendo fundamental para a identificação e
mapeamento de materiais. Frameworks de classificação baseados em características
espectral-espaciais superam métodos que utilizam apenas informações espectrais.
Sistemas embarcados para detecção de plásticos marinhos demonstram precisão
próxima a 90\% utilizando abordagens híbridas \cite{Alboody2023}, enquanto
implementações FPGA de SVM alcançam 99,7\% de precisão com latência de 0,1 ms por
pixel \cite{Martins2019}.

Entre os métodos clássicos de classificação, destacam-se as Support Vector
Machines (SVM), classificadores robustos que demonstram excelente desempenho em
dados hiperespectrais, com implementações otimizadas conseguindo precisões
superiores a 95\% em datasets de referência. A eficácia do SVM está relacionada à
sua capacidade de lidar com alta dimensionalidade e pequenos conjuntos de
treinamento. Implementações em FPGA conseguem classificação em tempo real com
latência de 0,1 ms por pixel e precisão de 99,7\% em datasets AVIRIS
\cite{Martins2019}.

Random Forest, métodos ensemble que combinam múltiplas árvores de decisão,
oferecem boa interpretabilidade e robustez a outliers, com tempos de treinamento
significativamente menores que SVM. O Spectral Angle Mapper (SAM) realiza a
classificação baseada em similaridade angular entre espectros, sendo
computacionalmente eficiente para implementações embarcadas e menos sensível a
variações de iluminação.

\section{Arquiteturas Heterogêneas de Processamento}
\label{sec:arquiteturas_heterogeneas}

As arquiteturas heterogêneas combinam diferentes tipos de processadores para
otimizar simultaneamente desempenho, consumo energético e latência. Essa
abordagem é particularmente relevante para o processamento hiperespectral
embarcado, onde diferentes estágios do pipeline computacional se beneficiam de
tipos distintos de processadores.

\subsection{Fundamentos de Computação Heterogênea}
\label{subsec:fundamentos_heterogenea}

A computação heterogênea fundamenta-se no princípio de que diferentes tipos de
processadores possuem eficiências variáveis para diferentes tipos de cargas
computacionais. A integração visa aproveitar os pontos fortes únicos de cada
arquitetura, melhorando a eficiência computacional e o desempenho em vários
domínios. \textcite{Vaithianathan2025} e \textcite{SahovicQuinten2025} destacam
a especialização de processadores, onde cada tipo (CPU, GPU, FPGA, DSP)
apresenta vantagens específicas para determinados tipos de operações. FPGAs, por
exemplo, destacam-se na aquisição rápida de dados, CPUs na coordenação de
sistema, e GPUs em computação massivamente paralela.

\begin{figure}[htbp]
    \centering
    \includegraphics[width=0.9\textwidth]{exemplo_imagem_hiperspectral_aviris}
    \caption{Exemplo de imagem hiperespectral AVIRIS capturada sobre São João D'Aliança, 
             GO, mostrando dados de vapor d'água em diferentes bandas espectrais. 
             A imagem ilustra a complexidade dos dados hiperespectrais que devem ser 
             processados por arquiteturas heterogêneas em sistemas embarcados para 
             aplicações de sensoriamento remoto em tempo real. 
             Fonte: \cite{Baptista2019}.}
    \label{fig:exemplo-imagem-hiperespectral-aviris}
\end{figure}

O particionamento de algoritmos, que consiste na divisão do pipeline de
processamento em estágios adequados para cada tipo de processador, é crucial. A
\autoref{fig:exemplo-imagem-hiperespectral-aviris} ilustra a complexidade dos dados 
hiperespectrais que motivam o uso de arquiteturas heterogêneas. A alocação 
estratégica de tarefas de acordo com os pontos fortes de cada componente de 
hardware atende às demandas computacionais e de tempo real essenciais para 
diversas aplicações \cite{SahovicQuinten2025}. A otimização
multi-objetivo busca o equilíbrio simultâneo entre desempenho, consumo
energético, latência e precisão. A alocação dinâmica de tarefas para a unidade
de processamento mais adequada otimiza a utilização de recursos e minimiza a
latência \cite{Vaithianathan2025}.

Finalmente, o codesign HW/SW, uma metodologia integrada de desenvolvimento que
considera hardware e software simultaneamente, enfatiza o desenvolvimento
simultâneo de componentes de hardware e software para alcançar desempenho e
eficiência ideais \cite{Vaithianathan2025}.

\subsubsection{Tendências Futuras em Computação Heterogênea}
\label{subsubsec:tendencias_futuras_heterogenea}

\textcite{Vaithianathan2025} identifica direções emergentes que moldarão
o futuro da computação heterogênea, com implicações diretas para o
processamento hiperespectral embarcado. A integração de aceleradores
específicos para IA, como Neural Processing Units (NPUs) e
Application-Specific Integrated Circuits (ASICs), oferece eficiência
energética cinco a dez vezes superior para inferência de IA comparado a
GPUs genéricas. Google TPUv4 demonstra 420 TFLOPS em operações de machine
learning, enquanto NPUs móveis alcançam 15-45 TOPS com consumo inferior a 5
W.

A heterogeneidade hierárquica evolui para sistemas que utilizam múltiplas
arquiteturas não apenas dentro de um único nó, mas também através de sistemas
distribuídos. Essa abordagem permite a combinação de diferentes processadores
otimizados para tarefas específicas em múltiplas camadas do sistema.
Arquiteturas de memória unificada, como AMD Heterogeneous System Architecture
(HSA) e NVIDIA Grace Hopper (900 GB/s de largura de banda CPU-GPU), reduzem 40\%
a 60\% da latência de transferência de dados, simplificando significativamente o
modelo de programação.

Frameworks de programação unificados desenvolvem abstrações de alto nível que
simplificam a programação cross-platform, permitindo que desenvolvedores
escrevam código uma vez e executem em múltiplas configurações de hardware sem
modificações extensivas. A alocação autônoma de recursos, com sistemas
orientados por IA que analisam métricas de performance em tempo real e ajustam a
alocação de recursos dinamicamente, otimiza a utilização e minimiza a latência
sem intervenção manual.

\subsection{Arquiteturas CPU+GPU Integradas}
\label{subsec:arquiteturas_cpu_gpu}

A combinação de CPU e GPU em sistemas embarcados oferece flexibilidade e
desempenho para o processamento hiperespectral. Arquiteturas de memória
unificada facilitam o compartilhamento contínuo de dados entre CPUs e GPUs,
reduzindo overhead e melhorando a velocidade computacional
\cite{Vaithianathan2025}. Sistemas como NVIDIA Jetson integram CPUs ARM com GPUs
CUDA em soluções single-package. A tendência em direção a designs System on Chip
(SoC) permite a integração de múltiplos tipos de processadores dentro de um
único chip, aprimorando ainda mais o desempenho em ambientes compactos
\cite{Vaithianathan2025}.

A divisão de tarefas aloca a CPU para controle, pré-processamento sequencial e
classificação final, enquanto a GPU é utilizada para operações matriciais,
filtragem e redes neurais. O controlador DMA atua como ponte entre CPU e GPU,
facilitando a transferência contínua de dados entre a memória host e a memória
global da GPU \cite{Vaithianathan2025}. A comunicação eficiente é garantida por
memória compartilhada e interfaces de alta velocidade para minimizar o overhead
de transferência de dados. Tecnologias como AMD's Heterogeneous System
Architecture (HSA) exemplificam essa tendência ao permitir o acesso de memória
compartilhada através de várias unidades de processamento
\cite{Vaithianathan2025}.

Estudos demonstram throughput superior a 330 fps em compressão hiperespectral
com consumo inferior a 20 W. Embora GPUs forneçam forte capacidade paralela,
encontram ineficiências na transferência de dados \cite{SahovicQuinten2025}.

\subsection{Sistemas CPU+FPGA}
\label{subsec:sistemas_cpu_fpga}

A combinação CPU+FPGA oferece máxima flexibilidade para otimização de algoritmos
específicos. Implementações bem-sucedidas de aceleradores de hardware para
classificação de imagens hiperespectrais em tempo real utilizando essa abordagem
já foram demonstradas \cite{Martins2019}. A FPGA implementa algoritmos críticos
como seleção de bandas EMCR e classificação SVM otimizada. A metodologia EMCR
(Entropy and Mean Correlation Ratio) aumentou consideravelmente o poder de
predição, alcançando níveis de precisão superiores a 99\% \cite{Martins2019}.

Um pipeline de baixa latência garante processamento contínuo com latência
determinística, adequado para aplicações de tempo real. FPGAs garantem baixo
nível de latência comparadas a GPUs e CPUs, pois trabalham em bare metal sem
sistema operacional. Melhorias de até 43,5 vezes no consumo energético são
alcançadas através de codesign otimizado. A natureza reconfigurável dos FPGAs
permite a maximização da otimização adaptada a aplicações específicas. Ao
contrário de CPUs e GPUs, que são arquiteturas fixas programáveis por software,
FPGAs são reconfiguráveis com engines computacionais definidos pelo usuário.
